
%\needspace{4\baselineskip}
\color{uniblue}{\subsection{Контроль отсутствия напряжения до ввода (КОНвв)}}
\color{black}

\begin{enumerate}[label=\arabic{section}.\arabic{subsection}.\arabic*, labelsep=4pt, leftmargin=0pt, itemindent=57pt]

\item
В устройстве контроль отсутствия напряжения до ввода представлен функцией «КОНвв». Функция используется для АПВ, а также для оперативного включения.

\item
Пусковые  условия  функции  формируются  без  выдержки  времени  при  снижении  всех  линейных напряжений до ввода ниже уставки и одновременном выполнении следующих условий, если уставка SGF2 выставлена в положение «Измеренный»:
\begin{itemize}
    \item функция введена в работу;
    \item минимальное из подведенных линейных напряжений, измеренных до ввода ниже уставки;
    \item отсутствует срабатывание ИО напряжения обратной последовательности;
    \item неактивный сигнал срабатывания БННвв.
\end{itemize}

\item
Предусмотрена блокировка функции при срабатывании ИО обратной последовательности напряжения выше уставки и наличии сигнала срабатывания БННвв.

\item
Если  уставка  SGF2 «Реж\_пуска»  выставлена  в  положение  «Внешний»,  работа  функции  осуществляется  при наличии  активного  внешнего  дискретного  сигнала  «Отсутствие  Uвв»  без  учета  срабатывания измерительных органов. 

\item
Ввод функции в работу на этапе параметрирования устройства осуществляется программным переключателем SGF1 «Ввод\_функции». Информация о введенной функции отображается выходным сигналом «КОНвв: Ввод».

\item
Функция может быть выведена из работы оперативно путем активации сигнала «ОВ КОНвв», что характеризуется активным состоянием сигнала «КОНвв: Оперативный вывод» на выходе функции.

\item
Алгоритм работы функции выполнен в соответствии с рисунком \ref{konvv:img1}. В таблице \ref{konvv:tbl1} приведены параметры, необходимые для настройки функции.

\vspace{1mm}
\begin{figure}[H]
\centering
\includegraphics[width=0.9\textwidth,height=0.9\textheight,keepaspectratio]{\fbpath/307. КОНвв/_latex/img/konvv.pdf}
\captionof{figure}{Функциональная схема логики КОНвв}\label{konvv:img1}
\end{figure}

\end{enumerate}
